\subsection{Riesgos y contingencias}

Dentro del proyecto pueden llegar a suscitarse situaciones en las que podamos tener retrasos o algunos fallos, a esto lo conocemos como riesgos del proyecto, para evitar que sucedan estas situaciones es necesario conocerlas y tenerlas presentes para tener cuidado con ello.

Para poder medir el riesgo como tal y no solo definirlo, nos basamos en 3 puntos importantes, la probabilidad de ocurrencia, el impacto que tendría y finalmente por medio de la multiplicación de estos 2, medir la magnitud del riesgo. La probabilidad de ocurrencia tiene una escala de valores de 0 a 100, medido en base a nuestras creencias. Para el impacto seguimos con una escala de valores enteros de 1 a 5 que representan lo siguiente:

\begin{itemize}
  \item \textbf{1:} Se considera insignificante, no afectaría en los resultados ni en fechas de entrega.
  \item \textbf{2:} Un nivel menor, se tendría que volver a trabajar pero no es tan grande como para modificar las fechas de entrega.
  \item \textbf{3:} Este valor se considera moderado ya que afectaría en la calidad de nuestro trabajo final, lo que seguramente resultaría en un retrabajo mayor al punto anterior.
  \item \textbf{4:} Ahora este valor se toma en cuenta como un conflicto mayor, esto comprometería completamente algunos de los objetivos del proyecto planteados anteriormente.
  \item \textbf{5:} Finalmente llegamos al nivel crítico, en donde se compromete el proyecto completo.
\end{itemize}

Por otro lado, para la magnitud se tuvieron en cuenta los siguientes puntos para determinar si llega a ser algo que se tenga que mitigar o no es necesario.

\begin{itemize}
  \item \textbf{$>$3.0:} Nivel de magnitud alto, lo que implicaría que es obligatorio mitigar ese riesgo.
  \item \textbf{1.5 a 3.0:} Nivel medio, en donde se tiene que mitigar de igual forma y tener un plan de contingencia, pero solamente si implica un bajo costo monetario y temporal.
  \item \textbf{$<$1.5:} Si es un valor menor a 1.5 se acepta el riesgo aunque se monitorea, ya que se consideraría de magnitud baja.
\end{itemize}

\begin{table}[H]
  \centering
  \small
  \begin{tabularx}{\textwidth}{>{\raggedright\arraybackslash}p{0.15\textwidth}>{\raggedright\arraybackslash}p{0.12\textwidth}>{\raggedright\arraybackslash}p{0.08\textwidth}>{\raggedright\arraybackslash}p{0.09\textwidth}>{\raggedright\arraybackslash}X>{\raggedright\arraybackslash}X}
    \toprule
    Riesgo & Probabi\-lidad de ocurrencia & Impacto & Magnitud & Plan de mitigación & Plan de contingencia \\
    \midrule
    Dataset chico o desbalanceado & 85\% & 5 & 4.25 & Definir un mínimo de vacas a las que tomarles foto, además de decidir cuántas fotos tomar por vaca. & Utilizar algún modelo preentrenado con capas congeladas si es posible. \\
    \addlinespace
    Calidad irregular en las fotos & 80\% & 5 & 4.0 & Definir una misma vista, distancia, altura, etc. para las diferentes vacas. Además de descartar fotos borrosas. & Usar data augmentation y volver a eliminar fotos irregulares. \\
    \addlinespace
    Data leakage por no separar por vacas & 80\% & 4 & 3.20 & Separar train, validation y test por vaca y no solo al azar. & Rehacer la separación por vaca y reentrenar. \\
    \addlinespace
    El modelo aprende cosas que no son de la vaca en sí & 75\% & 4 & 3.0 & Se puede recortar o segmentar a la vaca, además de poder implementar un bounding box a la imagen. & Segmentar y enmascarar el fondo de la imagen. \\
    \addlinespace
    Restricciones técnicas & 60\% & 3 & 1.80 & Utilizar un iPhone de al menos un integrante del equipo, además de conseguir una GPU para el entrenamiento. & Si no se puede conseguir una cámara con 3D, usar RGB. \\
    \bottomrule
  \end{tabularx}
\end{table}
